\section{协作式对话}
\label{sec:nlp-dialogue-collaboration}

用户说“把通知写得简单些”，并没有完整规定交付物。“简单”可能指篇幅、阅读难度或语气；直接改短有时合适，有时删掉了用户最关心的条件。协作式对话与用户共同确定并完成一件产物，挑战有三项：要求含混时，决定是否询问、问什么；多个可行方案难以抽象描述时，让用户在看得见的候选中选择，并把选择正确地记入要求；反馈只说“不满意”而不说怎么改，修改一处又可能牵动别处。对应的方法谱系是：目标澄清（规则询问、信息增益、学习的澄清问题排序）→候选讨论（差异候选构造、成对偏好排序）→互动修订（用户反馈条件编辑、模型自反馈修订）。生成与编辑的算法在第\ref{chap:nlp-text-generation}章，本节研究交流怎样决定要生成什么，以及产物是否落实了当前要求。

\subsection{目标澄清}
\label{sec:nlp-dialogue-clarification}

澄清的输入是含混请求、已有材料和已确认条件，输出是值得优先询问的问题。对“简单些”，可以直接问希望多短，可以给两个版本让用户比较，也可以先按默认理解交稿；选哪一种取决于误解的代价、问题是否好回答和用户愿意投入多少。已确定的活动事实不再询问，尚未确定的目标也不能由默认候选冒充用户确认。下面先给可解释的规则与信息增益基线，再介绍从数据中学习问题价值的方法。

\subsubsection{神经EVPI澄清问题排序}
\label{sec:nlp-dialogue-evpi}

规则和信息增益都需要人写出可能的目标与回答，换一个领域就要重写。Rao与Daumé III改为从论坛数据中学习：一个澄清问题有价值，是因为它的回答补上了原帖缺的信息\scite{nlp-dialogue-rao2018}。他们借用决策论中的\term{完全信息价值}（expected value of perfect information, EVPI）：如果知道问题$q$的回答$a$，原帖会变得多有用；对可能的回答取期望，就得到问$q$的价值。这一方法对候选问题排序，不生成新问题。

设当前帖子为$p$，候选问题为$q_i$。从相似帖子中检索出的问答对$(q_j,a_j)$提供可能的回答。LSTM编码帖子和问题，前馈网络$F_{\mathrm{ans}}$预测一个回答表示；$\widehat{\bm{a}}_j$与$\widehat{\bm{q}}_j$分别是候选回答与问题的平均词向量。记距离$d_{i,j}=1-\cos(F_{\mathrm{ans}}(p,q_i),\widehat{\bm{a}}_j)$，排序分数为
\begin{equation}
\begin{aligned}
 w_{i,j}&=\exp(-d_{i,j})\,
            \cos(\widehat{\bm{q}}_i,\widehat{\bm{q}}_j),\\
 S(q_i;p)&=\sum_j w_{i,j}U(p,q_j,a_j).
\end{aligned}
\label{eq:nlp-dialogue-evpi}
\end{equation}
$w_{i,j}$衡量“$a_j$像是$q_i$的回答”的程度，扮演回答概率的角色；$U$是效用网络，把帖子、问题与回答的LSTM表示送入前馈网络后经sigmoid输出，表示补上该回答后帖子的有用程度。原式未把$w_{i,j}$归一化为概率，这里保持原形式。训练有两个目标：回答表示损失让$F_{\mathrm{ans}}$的输出接近原配对回答；效用模型以原始的帖子—问题—回答三元组为正例、检索到的其他组合为负例做二分类。负例带噪声，因为别的问题也可能有用；监督来自论坛的交互记录，而非逐问题测得的真实效用。

对“把通知写得简单些”，取$q_1$为“需要限制篇幅吗”，$q_2$为“是否面向首次参加者”，候选回答$a_1$为“限制在80字左右”，$a_2$为“面向首次参加者”。设两个问题的交叉余弦为0.5、自身为1，距离矩阵为$\left[\begin{smallmatrix}0.2231&0.9163\\0.5108&0.5108\end{smallmatrix}\right]$，权重近似为$(0.8,0.2)$和$(0.3,0.6)$。给定效用$U_1=0.7$、$U_2=0.9$，得$S(q_1)=0.74$、$S(q_2)=0.75$，先问受众。展示问题后要等待用户的真实回答，不能把候选回答$a_2$写成已确认的事实。这一方法的限制也来自它的形式：候选里没有关于语气的问题就选不出它，效用分数也没有扣除用户回答的负担。

\subsubsection{规则询问与信息增益基线}
\label{sec:nlp-dialogue-information-gain}

规则基线先查必需条件：用途未定则问用途，两个硬要求冲突则先澄清冲突，其余采用声明的默认值并交付初稿。另一个可手算的基线把未知目标表示为有限个假设$G$，计算问题$q$的回答$A_q$能减少多少不确定性，再扣去询问成本$c(q)$：
\begin{equation}
 J(q)=H(G)-\sum_a p(a\mid q)H(G\mid A_q=a)-c(q),
 \label{eq:nlp-dialogue-information-gain}
\end{equation}
熵以2为底，前两项之差是回答带来的\term{信息增益}（information gain），即知道回答后目标熵的期望下降量。回答概率由假设先验与回答模型求得；若回答可能误解或拒答，需在$p(a\mid q,G)$中计入。把成本换算成与比特可比的单位是教学约定，不是普适标准。与上一小节相比，这个基线要求显式写出先验、回答和成本，换来的是每一步都可以复核。

规定四个互斥的教学目标：只缩短篇幅$g_1$，只降低难度$g_2$，只放松语气$g_3$，同时缩短并降低难度$g_4$，先验为$(0.4,0.3,0.1,0.2)$。问题$q_L$问“是否主要需要更短”，把$g_1,g_4$归入“是”，回答概率为$(0.6,0.4)$；$q_D$问“是否需要降低阅读难度”，把$g_2,g_4$归入“是”，概率为$(0.5,0.5)$。假设回答确定且诚实。先验熵为1.846比特。$q_L$的“是”后验为$(2/3,1/3)$，“否”后验为$(3/4,1/4)$，预期剩余熵$0.6\times0.918+0.4\times0.811=0.875$，信息增益0.971；$q_D$两组后验为$(0.6,0.4)$与$(0.8,0.2)$，预期剩余熵0.846，信息增益1。若$c(q_L)=0.10$、$c(q_D)=0.25$，净分数为0.871与0.750，优先问篇幅。

直接按最高先验目标$g_1$交稿，猜对概率为0.4；问$q_L$后各分支取最大后验目标，命中概率升到$0.4+0.3=0.7$。多问一轮的收益由此可见，但具体字数和难度等级仍未确定。熵减只衡量当前假设空间内的不确定性，可能与产物是否可用不一致；直接交付、给选项与先询问三种做法最终仍要按交付质量和用户投入比较。

\subsection{候选讨论}
\label{sec:nlp-dialogue-candidates}

用户常常说不清想要什么语气，却能看出哪个版本合适。候选讨论要让差异可辨，又不能拿一份正确、一份漏了事实的版本制造伪选择。实现上分两步：构造覆盖硬条件、彼此有实质差异的候选；把用户的选择转成对候选的偏好分数。用户点选某个版本或说“挺好”，只按当时问题的范围解释，不自动视为所有细节都已验收。

\subsubsection{约束覆盖与差异候选选择}
\label{sec:nlp-dialogue-diverse-candidates}

固定任务为发布当前活动简讯，已确认事实包括活动名、当前时间地点、校内参加条件和校外严格截止条件。生成器先以这些条件产生候选，规则检查字段覆盖；只有通过硬条件的候选才进入篇幅、语气与组织方式的比较，避免用多样性弥补事实缺失。例如版本甲为“机器学习读书会将于2025年4月18日（周五）14:00在海棠楼201举行。校内师生可直接参加，校外来宾须在4月16日18:00前提交报名信息。”版本乙为“机器学习读书会：2025年4月18日14:00，海棠楼201。校内师生可直接参加；校外来宾须在4月16日18:00前提交报名信息。”两者事实相同，差异可以说明为完整叙述与紧凑简讯。

候选较多时，用第\ref{chap:nlp-text-generation}章的最大边际相关性（MMR）逐个选入：设相关性为$R(y)$、已选集合为$S$，
\begin{equation}
 s(y)=\lambda R(y)-(1-\lambda)\max_{z\in S}\operatorname{sim}(y,z).
 \label{eq:nlp-dialogue-mmr}
\end{equation}
已选甲、$\lambda=0.7$时，乙的相关性0.9、与甲相似度0.6，得分0.45；几乎同文的丙为0.95与0.95，得分0.38，于是选乙。展示数量要计入用户的阅读负担，比较不同展示方案时应固定生成预算与展示条数。

\subsubsection{成对偏好下的候选排序}
\label{sec:nlp-dialogue-bradley-terry}

用户在多轮中反复二选一，积累的是成对比较。把它们变成候选排序，用的是第\ref{sec:nlp-align-bradley-terry}小节的Bradley--Terry模型：给每个候选一个分数$s_i$，选$i$胜过$j$的概率为$\sigma(s_i-s_j)$，用比较次数的负对数似然估计分数；模型的推导、零点不可识别与全胜发散在那里已经讲过。用在协作对话中，要处理的是三件事。

第一是读出分差。同一用户在相同目标下比较甲乙四次，乙胜3次、甲胜1次，固定$s_{\text{甲}}=0$，最大似然估计为$s_{\text{乙}}=\log3\approx1.10$，即乙的胜负比是3比1。这一估计把四次选择当作独立同分布，而顺序、记忆和疲劳都会影响重复询问，所以独立性是假设。第二是平局与带修改的反馈。“两个都行”“选乙但把日期恢复”不是一次单纯的乙胜：平局按事先约定剔除或用支持平局的扩展模型，修改意见作为独立反馈保留，交给下一小节的编辑。第三是作用范围。分数描述的是某个问题、某个用户条件下的相对选择，不证明乙的每条事实都对；聚合其他用户偏好得到的分数，也不能覆盖当前用户已经明确的要求。

用户说“第二个更短，但保留地点”，共同状态就更新为采用乙的紧凑形式，并把“地点必须保留”记为已确认条件。乙本来含地点，无需为显示“已修改”再改一次；是否进一步缩短、是否省略年份仍未确认。

\subsection{互动修订}
\label{sec:nlp-dialogue-revision}

选定候选后，交流的结果要落到产物上。系统保存稿件版本号、已确认条件和未决问题，把每条反馈绑定到它所针对的版本，避免用户评论旧稿时误改新稿。挑战在于反馈常只指出不满意之处，且修改一处可能要同步另一处。两条路线的区别在反馈从哪里来：用户反馈条件编辑依赖用户给出的意见，Self-Refine让模型自己生成反馈。两者都要先确定作用范围再编辑，以实际差异和验收判断是否完成，而不是以回复了一句“已修改”为准。

\subsubsection{用户反馈条件编辑}
\label{sec:nlp-dialogue-feedback-edit}

用户反馈条件编辑把（原稿$y$，反馈$f$）作为条件，生成修订稿$y'$。设$x$为事实材料与已确认条件，训练数据为原稿—反馈—修订稿三元组，目标是修订稿的词元负对数似然：
\begin{equation}
 L_{\mathrm{edit}}=-\log p_{\theta}(y'\mid x,y,f)
  =-\sum_{i=1}^{|y'|}\log p_{\theta}(y'_i\mid x,y,f,y'_{<i}).
 \label{eq:nlp-dialogue-feedback-edit}
\end{equation}
模型结构可以是把$x$、$y$、$f$拼接为源序列的编码器—解码器，也可以是第\ref{sec:nlp-generation-editing}节那样预测保留、删除、替换操作的局部编辑器；提示式方法则在固定接口下由大模型直接改写。与第\ref{chap:nlp-text-generation}章的一般编辑相比，新增的是反馈$f$：它只说明要改什么，没有给出改后全文。

训练目标只要求输出接近参考修订稿，不检查改动是否落在反馈要求的范围内，因此需要单独的编辑检查。用Levenshtein对齐求出原稿到修订稿的改动位置集合$E(y,y')$，再由反馈与当前状态确定必须改动的位置集合$R(f)$，定义
\begin{equation}
 \mathrm{NMR}=\frac{\lvert E(y,y')\cap R(f)\rvert}{\lvert R(f)\rvert},
 \qquad
 \mathrm{EXR}=\frac{\lvert E(y,y')\setminus R(f)\rvert}{\lvert E(y,y')\rvert}.
 \label{eq:nlp-dialogue-edit-rates}
\end{equation}
NMR是必要修改完成率，EXR是改动中未经授权的比例；两者都是本书为检查编辑范围给出的教学指标。$R(f)$必须由反馈与已确认事实确定，例如“改日期”确定的是所有表示当前活动日期的位置，而不是所有日期字样。

设教学旧稿有两处“4月17日”，用户说“改成4月18日，地点不变”，依据是共同通知中已确认的改期。用压缩词元序列表示：旧稿为“读书会／4月17日／海棠楼201／请于／4月17日／到场”，目标为“读书会／4月18日／海棠楼201／请于／4月18日／到场”，$R(f)$为第2、5个位置。若修订稿只替换了第2个位置，NMR为$1/2$；若两处都替换但把地点改成另一栋楼，NMR为1而EXR为$1/3$。若原稿另有一句“原定4月17日”，它表示旧计划，不属于$R(f)$。有歧义的“改到周五”需先绑定日期属性和年份，不能改报名截止；反馈与权威事实冲突时，要先澄清是在编辑假设通知还是修订真实材料。

\subsubsection{Self-Refine辅助自反馈修订}
\label{sec:nlp-dialogue-self-refine}

用户反馈条件编辑需要用户逐条指出问题。Self-Refine让同一个语言模型依次承担初稿生成、反馈和修订，不训练专门的反馈器，也不更新参数；三种角色分别由三类提示中的说明与少量示例指定\scite{nlp-dialogue-madaan2023}。设请求与材料为$x$，初稿为$y_0$，第$t$轮反馈为$f_t$：
\begin{equation}
\begin{aligned}
 y_0&=M(p_{\mathrm{gen}},x),\\
 f_t&=M(p_{\mathrm{feedback}},x,y_t),\\
 y_{t+1}&=M(p_{\mathrm{refine}},x,y_0,f_0,\ldots,y_t,f_t).
\end{aligned}
\label{eq:nlp-dialogue-self-refine}
\end{equation}
修订提示保留此前全部输出与反馈，循环到固定轮数或反馈中出现停止指示为止。停止指示本身也是模型输出。

假设初稿把要求写成“校外来宾在4月16日前报名”。反馈提示要求“对照给定材料，定位被删去的时间与限定词，只指出可执行的修改”，得到“补回18:00及严格的前，其他字段保持不变”；修订后为“校外来宾须在4月16日18:00前提交报名信息”。再用规则检查截止时刻、严格不等号、主体与动作，才确认条件未变。这一轮共有生成、反馈、修订三次调用。

与用户反馈条件编辑相比，Self-Refine减少了对用户逐条指出问题的依赖，代价是多次模型调用，而且模型可能察觉不到自己的错误，或在反馈中重复同一错误。自反馈的优先级低于用户已确认的要求：若它建议“删除地点以便更短”，而地点已被确认必须保留，编辑器应拒绝这条建议。评价时分别记录用户轮数与自循环调用数，收益由最终产物的独立检查衡量。

表\ref{tab:nlp-dialogue-collaboration-methods}汇总本节方法。

\begin{table*}[htbp]
  \centering
  \small
  \setlength{\tabcolsep}{3pt}
  \renewcommand{\tabularxcolumn}[1]{>{\raggedright\arraybackslash}p{#1}}
  \begin{tabularx}{\linewidth}{@{}>{\raggedright\arraybackslash}p{28mm}XXX@{}}
    \toprule
    方法 & 解决的问题 & 核心机制 & 新代价\\
    \midrule
    规则询问 & 缺少必需条件或硬要求冲突 & 按条件清单询问，其余用声明的默认值 & 规则覆盖有限\\
    \addlinespace
    信息增益 & 不知先问哪个问题 & 对有限目标假设计算熵减并扣除成本 & 需写出先验、回答模型与成本单位\\
    \addlinespace
    神经EVPI排序 & 人写假设难以迁移 & 从论坛问答学习回答表示与效用，按期望效用排序 & 只能排序已有候选；不计用户负担\\
    \addlinespace
    差异候选选择 & 用户难以抽象描述偏好 & 硬条件过滤后用MMR选有差异的候选 & 用户阅读负担；展示数量需控制\\
    \addlinespace
    成对偏好排序 & 多次选择需汇总 & Bradley--Terry把成对选择转为分数 & 独立性假设；平局与修改需另记\\
    \addlinespace
    用户反馈条件编辑 & 意见需落到具体位置 & 以原稿与反馈为条件生成修订稿，再检查改动范围 & 需要三元组数据或固定接口；依赖用户逐条反馈\\
    \addlinespace
    Self-Refine & 用户逐条反馈负担大 & 同一模型循环生成、反馈、修订 & 调用次数增加；自评不能替代核验\\
    \bottomrule
  \end{tabularx}
  \caption{协作式对话的方法对照。前三行处理目标澄清，中间两行处理候选讨论，后两行处理互动修订；后一组方法都以前一组确定的要求为条件。}
  \label{tab:nlp-dialogue-collaboration-methods}
\end{table*}
